\documentclass[11pt]{article}
\usepackage{fontspec}
\setmainfont{Times New Roman}
\setsansfont{Arial}
\setmonofont{Consolas}
\usepackage{amsmath,amssymb}
\usepackage{geometry}
\usepackage{microtype}
\geometry{a4paper,margin=3cm}
\setlength{\parindent}{1.25em}
\setlength{\parskip}{0pt}
\emergencystretch=30em
\allowdisplaybreaks
\raggedbottom
\providecommand{\mM}{\mathfrak M}
\providecommand{\mN}{\mathfrak N}
\providecommand{\mA}{\mathfrak A}
\providecommand{\mB}{\mathfrak B}
\providecommand{\mP}{\mathfrak P}
\providecommand{\mD}{\mathfrak D}
\providecommand{\mS}{\mathfrak S}
\providecommand{\ma}{\mathfrak a}
\providecommand{\mb}{\mathfrak b}

\begin{document}

% Унит: Папер 17 Сецтион 12 анд енд маттер. Дирецт Интерславиц Латин транслатион фром те Герман финал-аудитед соурце слице.

\setcounter{footnote}{23}
\subsection*{\S~12. Примеры.}

1. Како први примјер разсмотримо обычно линеарно диференциално уравненье 2-го реда, чије коефициенты сут представјене како симетричне функције интегралов. Област рационалности \(P\) ту состоји из всех рационалных функциј од интегралов и јих дериватов, при чем се ограничујемо на довольно малу околину регуларној точкы диференциалного уравненьа. Модул \(\mathfrak M\) определьајемо како целокупност всех диференциалных изразов, кторе имајут \(y_1\) и \(y_2\) за интегралы; за ньих зново пишемо полиномы в једнеј неизвестној \(\xi\). Базис модула тогда јест једноелементны, именно даны полиномом
\[
 M=
 \begin{vmatrix} y_1&y_2\\ y_1'&y_2'\end{vmatrix}\xi^2
 -\begin{vmatrix} y_1&y_2\\ y_1''&y_2''\end{vmatrix}\xi
 +\begin{vmatrix} y_1'&y_2'\\ y_1''&y_2''\end{vmatrix},
\]
кторы јест једнозначно определьен до фактора из \(P\). Ныне
\begin{equation}
 \mathfrak M=[\mathfrak M_1,
\mathfrak M_2],
\tag{41}
\end{equation}
где \(\mathfrak M_i\) состоји из всех диференциалных изразов, кторе имајут интеграл \(y_i\), и јего базис јест, до фактора из \(P\), определьен чрез
\[
 M_i=y_i\xi-y_i'\qquad(i=1,2).
\]
В самом делу, \(\mathfrak M\) јест делимы чрез \(\mathfrak M_1\) и \(\mathfrak M_2\), бо изчезаје за интегралы \(y_1\) и \(y_2\), теды јест делимы чрез \([\mathfrak M_1,
\mathfrak M_2]\); и понеже ред базисного полинома од \([\mathfrak M_1,
\mathfrak M_2]\) мора быти најмене 2, имајемо \(\mathfrak M=[\mathfrak M_1,
\mathfrak M_2]\). Далеј, \(\mathfrak M_1\) и \(\mathfrak M_2\) сут взаимно просте, понеже
\begin{equation}
 \frac{y_2}{D}(y_1\xi-y_1')-\frac{y_1}{D}(y_2\xi-y_2')=1,
 \qquad
 D=D(y)=\begin{vmatrix}y_1&y_2\\ y_1'&y_2'\end{vmatrix}.
\tag{42}
\end{equation}
Кроме того они сут просте модулы, бо јих група остатков имаје ред 1.

Але \(\mathfrak M\) мимо \(y_1\) и \(y_2\) имаје такоже все интегралы
\[
 z_1=\alpha_{11}y_1+\alpha_{12}y_2,
 \qquad
 z_2=\alpha_{21}y_1+\alpha_{22}y_2,
 \qquad
 A=|\alpha_{ik}|\ne0,
\]
и зато јест идентичны с најменшим обчим многократным двух взаимно простых модулов \(\mathfrak N_1\) и \(\mathfrak N_2\), кторе имајут соответне интегралы \(z_1\) и \(z_2\); јих базисне полиномы сут даны чрез
\[
 N_i=z_i\xi-z_i'=(\alpha_{i1}y_1+\alpha_{i2}y_2)\xi-(\alpha_{i1}y_1'+\alpha_{i2}y_2').
\]
Тако \(\mathfrak M\) допушча бесконечно много различних представјень како најмене обче многократно взаимно простых простых модулов, и по Теорему VII оба модулы \(\mathfrak N_1\) и \(\mathfrak N_2\) меджу собоју и со всеми \(\mathfrak M_1\) и \(\mathfrak M_2\) сут једнакого вида. В самом делу, групы остатков всех тых модулов имајут ред 1, теды сут изоморфне.

Ныне запишемо адитивны разклад групы остатков \(\mathfrak G=\mathfrak A_1+\mathfrak A_2\), кторы соответује представјеньу (41). За то идемо, како в \S~4, од релације (42), ктора указује взаимну простост \(\mathfrak M_1\) и \(\mathfrak M_2\). За јединице \(a_1\) и \(a_2\) доставајемо класы остатков, породжене полиномами
\[
 -\frac{y_1}{D}(y_2\xi-y_2')
 \quad\text{и}\quad
 \frac{y_2}{D}(y_1\xi-y_1')
\]
соответне. Туто даје
\[
 \mathfrak M_1=\left(\mathfrak M,\frac{y_2}{D}(y_1\xi-y_1')\right),
 \qquad
 \mathfrak M_2=\left(\mathfrak M,-\frac{y_1}{D}(y_2\xi-y_2')\right).
\]
Ако ныне подобно творимо јединице \(b_1\) и \(b_2\), кторе соответујут разкладу \(\mathfrak M=[\mathfrak N_1,
\mathfrak N_2]\), заменьајучи \(y_i\) с \(z_i\), доставајемо
\[
 -\frac{z_1}{D(z)}(z_2\xi-z_2')
 =-\frac{\alpha_{11}y_1+\alpha_{12}y_2}{AD(y)}
 ((\alpha_{21}y_1+\alpha_{22}y_2)\xi-(\alpha_{21}y_1'+\alpha_{22}y_2')),
\]
\[
 \frac{z_2}{D(z)}(z_1\xi-z_1')
 =\frac{\alpha_{21}y_1+\alpha_{22}y_2}{AD(y)}
 ((\alpha_{11}y_1+\alpha_{12}y_2)\xi-(\alpha_{11}y_1'+\alpha_{12}y_2')),
\]
теды
\begin{equation}
\left\{
\begin{aligned}
 b_1&=\frac{\alpha_{11}y_1+\alpha_{12}y_2}{y_1}
 \left(\frac{\alpha_{22}}{A}\right)a_1
 +\frac{\alpha_{11}y_1+\alpha_{12}y_2}{y_2}
 \left(-\frac{\alpha_{21}}{A}\right)a_2
 =\frac{z_1}{y_1}\alpha^{11}a_1+\frac{z_1}{y_2}\alpha^{12}a_2,\\
 b_2&=\frac{\alpha_{21}y_1+\alpha_{22}y_2}{y_1}
 \left(-\frac{\alpha_{12}}{A}\right)a_1
 +\frac{\alpha_{21}y_1+\alpha_{22}y_2}{y_2}
 \left(\frac{\alpha_{11}}{A}\right)a_2
 =\frac{z_2}{y_1}\alpha^{21}a_1+\frac{z_2}{y_2}\alpha^{22}a_2.
\end{aligned}\right.
\tag{43}
\end{equation}
Ту \((\alpha^{ik})\) јест обратна матрица к матрици \((\alpha_{ik})\).

Обратно, чрез замену \(z_i\) с \(y_i\) доставаје се
\begin{equation}
\left\{
\begin{aligned}
 a_1&=\frac{y_1}{z_1}\alpha_{11}b_1+\frac{y_1}{z_2}\alpha_{12}b_2,\\
 a_2&=\frac{y_2}{z_1}\alpha_{21}b_1+\frac{y_2}{z_2}\alpha_{22}b_2.
\end{aligned}\right.
\tag{44}
\end{equation}
Из тых формул видна јест изоморфија груп \(\mathfrak A\) и \(\mathfrak B\). По нашим обчим разсмотрјеньам, при \(a_1b_\sigma\ne0\) соответно придруженье јест \(a_1\sim a_1b_\sigma\), теды по (44)
\[
 a_1\sim\frac{y_1}{z_\sigma}\alpha_{1\sigma}b_\sigma,
 \qquad\text{ако }\alpha_{1\sigma}\ne0,
\]
и по (43) подобно
\[
 b_\sigma\sim\frac{z_\sigma}{y_2}\alpha^{\sigma2}a_2,
 \qquad\text{ако }\alpha^{\sigma2}\ne0,
\]
теды
\[
 \frac{y_1}{z_\sigma}\alpha_{1\sigma}b_\sigma
 \sim
 \frac{y_1}{z_\sigma}\alpha_{1\sigma}\frac{z_\sigma}{y_2}\alpha^{\sigma2}a_2
 =\frac{y_1}{y_2}\alpha_{1\sigma}\alpha^{\sigma2}a_2.
\]
Тако јест
\[
 a_1\sim\alpha_{1\sigma}\alpha^{\sigma2}\frac{y_1}{y_2}a_2
 \quad\text{за }\alpha_{1\sigma}\alpha^{\sigma2}\ne0,
 \qquad
 \frac{1}{\alpha_{1\sigma}\alpha^{\sigma2}}\frac{y_2}{y_1}a_1\sim a_2,
\]
што даваје обратно придруженье. Ако \((\alpha_{ik})\) не јест диагонална матрица, то условје \(\alpha_{1\sigma}\alpha^{\sigma2}\ne0\) јест всегда изполньено за некоје \(\sigma\). Случај диагоналној матрице треба изкльучити, бо тогда модулы сут идентичне. Ако, на примјер, само \(\alpha_{12}=0\), тогда \(\mathfrak A_1=\mathfrak B_1\), але не \(\mathfrak A_2=\mathfrak B_2\). Из \(\mathfrak A_1+\mathfrak A_2=\mathfrak B_1+\mathfrak B_2\), \(\mathfrak A_1=\mathfrak B_1\), теды не следује \(\mathfrak A_2=\mathfrak B_2\) (срав. опазку 12). Три модулы \(\mathfrak M_1\), \(\mathfrak M_2\) и \(\mathfrak N_2\) сут тогда попарно взаимно просте, але не творјут тотално взаимно просту систему, бо \([\mathfrak M_1,
\mathfrak M_2]\) јест делимы чрез \(\mathfrak N_2\) (срав. опазку 15).

Далеј, ако положимо
\[
 t_{12}=\alpha_{1\sigma}\alpha^{\sigma2}\frac{y_1}{y_2}a_2,
 \qquad
 t_{21}=\frac{1}{\alpha_{1\sigma}\alpha^{\sigma2}}\frac{y_2}{y_1}a_1,
\]
то релације (35), како се потврджује изчисленьем, сут изполньене.

2. Примјер бесконечној групы даваје всакы модул од више неж једнеј пременној, кторы имаје базис само из једного полинома. Мало сложнејши јест следујучи примјер, кторы имаје показати, же бесконечне групы могут наступати такоже при многочленных модулах, т. ј. при такых, чиј модулны базис обхвача више неж једен полином.\footnote{Срав. примјер Ландауа, поданы од Блумберга (цит. дело, с. 51--52). Ту и в следујучем означујемо независне пременне с \(x\) и \(y\).}
\[
 \mathfrak M=[(\xi+x\eta),(\xi+1)]=[(P),(Q)]
 \qquad\left(\xi=\frac{\partial}{\partial x},\ \eta=\frac{\partial}{\partial y}\right).
\]
Правило множеньа нехај буде оно за диференциалне изразы. Група остатков ту ставаје бесконечна, бо групы остатков од \((\xi+x\eta)\) и \((\xi+1)\) сут обе бесконечне, доколе модулы \((\xi+x\eta)\) и \((\xi+1)\) сут взаимно просте понеже
\[
 1=(-x\xi-x^2\eta+1)Q+(x\xi+x-1)P.
\]
Ныне можно велми легко показати, же \(\mathfrak M\) не имаје полинома другого степена, положајучи
\[
 (u_1\xi+u_2\eta+u_3)(\xi+x\eta)=(v_1\xi+v_2\eta+v_3)(\xi+1).
\]
Из того следује \(u_1=u_2=u_3=v_1=v_2=v_3=0\). Ако далье ишчемо полиномы третјего степена в \(\mathfrak M\), находимо два линеарно независне, на примјер
\[
 (\xi^2+x\xi\eta+\xi+(2+x)\eta)Q=(\xi^2+2\xi+1)P
\]
и
\[
 (\xi^2+2x\xi\eta+x^2\eta^2+\eta)Q
 =(\xi^2+x\xi\eta+\xi+(x-1)\eta)P.
\]
Ако бы ныне \(\mathfrak M\) был једноелементны модул, тогда оба бы мусили быти делимы чрез базисны полином; и понеже \(\mathfrak M\) не имаје полинома другого степена, мусили бы быти делимы чрез полином третјего степена, теды быти пропорционалне, что не јест случај. Многочленност модула \(\mathfrak M\) јест внутришньа причина неуспеха Ландауовых продуктных теоремов, кторе важут в случају једнеј пременној.

3. Дальши примјер показује, же бесконечне групы могут наступати такоже при простых модулах.\footnote{Наша дефиниција ``простого модула'' јест различна од того, что в комутативном случају называје се простым модулом, понеже там всегда обставајут делительи нижшеј многовидности, кроме когда многовидност \(=0\), теды група остатков јест конечна.} \(P\) нехај буде област всех рационалных функциј двух пременных \(x,y\). Модул \(\mathfrak M\) с базисом
\[
 \xi-\frac{y}{x}=\xi-a
\]
имаје бесконечну групу, бо она обсахује клас остатков всакого степена \(y\); с другој страны, \(\mathfrak M\) јест просты модул. В самом делу покажемо, же всакы делитель
\[
 \mathfrak M_1=(\xi-a,
F(\xi,\eta))
\]
ставаје јединичным модулом. Всакы полином \(F(\xi,\eta)\) можно намреч мод\((\xi-a)\) представити полиномом \(G(\eta)\):
\[
 F(\xi,\eta)=c_0\eta^\nu+c_1\eta^{\nu-1}+\cdots+c_\nu\quad(\mathfrak M).
\]
Без ограниченьа обчности нехај \(c_0=1\). Ныне по закону множеньа \(\xi\eta^\nu-\eta^\nu\xi=0\); теды, понеже
\[
 \xi\equiv a(\mathfrak M_1),
 \qquad
 \eta^\nu\equiv F_1(\mathfrak M_1),
\]
где положено \(F_1=-(c_1\eta^{\nu-1}+\cdots+c_\nu)\), доставајемо
\[
 \xi F_1-\eta^\nu a\equiv0(\mathfrak M_1).
\]
Ако развијемо
\[
 \eta^\nu a=a\eta^\nu+\nu\frac{\partial a}{\partial y}\eta^{\nu-1}+\cdots,
\]
и потом ту зново заменимо \(\eta^\nu\) чрез \(F_1\), доставајемо
\[
 \xi(c_1\eta^{\nu-1}+\cdots+c_\nu)-a(c_1\eta^{\nu-1}+\cdots+c_\nu)
 +\nu\frac{\partial a}{\partial y}\eta^{\nu-1}+\cdots\equiv0(\mathfrak M_1),
\]
или
\[
 c_1\eta^{\nu-1}a+\frac{\partial c_1}{\partial x}\eta^{\nu-1}+\cdots
 -ac_1\eta^{\nu-1}+\cdots+
u\frac{\partial a}{\partial y}\eta^{\nu-1}+\cdots
 \equiv0(\mathfrak M_1),
\]
или наконец
\[
 \left(\frac{\partial c_1}{\partial x}+\nu\frac{\partial a}{\partial y}\right)\eta^{\nu-1}+\cdots
 \equiv0(\mathfrak M_1).
\]
На левој стране стоји ныне полином в \(\eta\) точного степена \(\nu-1\), кторы не изчезаје идентично; бо највышши коефициент
\[
 \frac{\partial c_1}{\partial x}+\nu\frac{\partial a}{\partial y}
 =\frac{\partial c_1}{\partial x}+\nu\frac{1}{x}
 \quad\text{јест }\ne0,
\]
понеже иначе \(c_1=-\nu\log x+\varphi(y)\) не был бы рационалны. Тако можно поступ наставити и он води к не идентично изчезајучему полиному нултого степена; тым показано, же такоже јединица јест в \(\mathfrak M_1\).\footnote{Туто поступанье своджи се к доказаньу, же необходне условја интеграбилности за систем диференциалных уравнень не сут изполньене.}

4. Ако, с другој страны, адјунгујемо области рационалности такоже функцију \(\log x\), тогда \(\mathfrak M\) имаје делителье
\[
 \left(\xi-\frac{y}{x},\ \eta-\log x+\varphi(y)\right),
\]
где \(\varphi(y)\) пробегаје все рационалне функције од \(y\). Ти делительи сут вси различни меджу собоју; за всакого јест изполньено условје интеграбилности, и соответны интеграл јест
\[
 y\log x-\int\varphi(y)\,dy+c.
\]
Зато јединица не може быти в модулу, бо иначе нулта функција была бы једины можны интеграл. С другој страны, \(\xi\) и \(\eta\) сут гледе на модул конгруентне некојеј величине области; зато група остатков јест конечна и имаје ред 1. Далеј, из бесконечно многих делительев всака два сут взаимно проста, и даны модул \(\mathfrak M\) јест равен најменшему обчему многократному \(\mathfrak N\) всех тых делительев. В самом делу, он јест делимы чрез все делителье, теды такоже чрез \(\mathfrak N\). Ако бы ныне \(\mathfrak N\) был властны делитель од \(\mathfrak M\), тогда \(\mathfrak N\) мусил бы имети јешче најмене једен полином \(F\), рецимо степена \(\rho\), кторы зависи само од \(\eta\). Ако ныне допустимо, же \(\varphi(y)\) пробегаје функције \(y,
\ldots,y^{\rho+1}\), тогда меджу соответными интегралам
\[
 y\log x-\frac{y^{i+1}}{i+1}+c_i\qquad(i=1,
\ldots,
\rho+1)
\]
не обставаје линеарна релација с коефициентами независными од \(y\). Але диференциално уравненье \(\rho\)-того реда \(F=0\) не може имети \(\rho+1\) линеарно независных интегралов.

Модул \(\mathfrak M=\mathfrak N\) але не јест пополно редуцибилны.\footnote{Тым имајемо примјер за први случај Теорема X.} Бо иначе был бы најменше обче многократно конечно многих од тых простых модулов, и јего група остатков имела бы конечны ред; медјутым група остатков од \(\xi-\frac{y}{x}\) обсахује класы остатков всех степенов \(\eta\), теды јест бесконечна.

Готтинген, 1 августа 1919.

\begin{center}
(Пријето 4 августа 1919.)
\end{center}

\end{document}

